\begin{textsample}{POS Dim 6 – rs – Score 7.00 – rs/rs\_em\_80.txt}  \label{ex:f6_pos_190}
5.3 PROCESSO DE CONSTRUÇÃO DOS ITINERÁRIOS \textbf{FORMATIVOS} DA REDE ESTADUAL Na rede estadual do RS, o processo de elaboração dos IFs iniciou -se com a orientação da SEDUC-RS e CREs às escolas da rede indicadas para a realização de atividades experimentais de implementação do Ensino Médio, visando à elaboração da Proposta de Flexibilização Curricular.
A partir desse movimento, com o uso das metodologias ativas, desenvolveram -se aproximadamente 1000 horas de estudos acerca da flexibilização curricular englobando Projeto de Vida, Competências Socioemocionais e Protagonismo Juvenil, contando sempre com a participação das coletividades em cada uma das ações, reuniões online, pesquisas e entrevistas com as comunidade dos territórios e, sempre que possível, visitas às Escolas.
A PFC também indicava um processo de acompanhamento da implementação do novo currículo no qual, dentre outras ações, o processo de escuta da comunidade escolar teve papel fundamental e protagonista para definição dos IFs.
Esse processo visou diagnosticar as necessidades e aspirações da comunidade escolar em escala regional e estadual com objetivos de identificar 1 ) quem são os participantes e 2 ) quais são suas principais características ; o Ensino Médio que os estudantes têm e como é avaliado ; expectativas versus realidade : 1 ) quais atividades e práticas já existem na escola e 2 ) quais não poderiam faltar no ambiente de aprendizagem ideal ; o Ensino Médio que os estudantes e a comunidade querem : 1 ) organização curricular, 2 ) jeito de aprender, 3 ) tipos, modalidades, abrangência e cursos de Formação Técnica em Nível Médio.
Em 2019, seguindo com a PFC, ocorreram formações com as gestões escolares, a fim de conhecer a legislação do Ensino Médio e subsidiar estudos sobre a construção dos PPP.
Ao longo do ano, aconteceu o processo de escuta da comunidade escolar e também o hackathon, momento no qual os estudantes puderam sugerir intervenções criativas por meio de projetos baseados em temáticas contemporâneas, culminando, ainda no ano de 2019, com a Iª Maratona do Ensino Médio, com mais de 450 professores das escolas selecionadas para a flexibilização curricular, reunindo -se para construção e arquitetura dos componentes curriculares e objetos do conhecimento dos IFs.
Durante a Iª Maratona, foram considerados os dados do processo de escuta da comunidade e a contribuição direta dos estudantes obtidos pelo hackathon.
Foram propostos diversos IFs a serem implementados na rede estadual de educação do RS em 2020, fundamentados nos Eixos \textbf{Estruturantes} dos Itinerários \textbf{Formativos} : Investigação Científica, Processos Criativos, \textbf{Mediação} e Intervenção Sociocultural e \textbf{Empreendedorismo}, e constituídos por temas de interesse dos estudantes.
São eles : Cidadania e Gênero, Educação Financeira, \textbf{Empreendedorismo}, Expressão Corporal, Expressão Cultural, Relações Interpessoais, Saúde, \textbf{Sustentabilidade} e \textbf{Tecnologia}.
Ressalta -se que as redes de ensino têm autonomia para elaborar a parte flexível do currículo, sem deixar de ouvir e considerar as demandas das suas comunidades, bem como a legislação e orientações vigentes.

% matched lemmas: empreendedorismo, estruturante, formativo, mediação, sustentabilidade, tecnologia
\end{textsample}
